% sec-04: the corporate reserve.  Figure 3 (d2-type grid with a measure strip),
% Table 4.  No dollar amount appears: the instruction sizes the reserve as a
% share, and a figure or table that invented a balance would state a number the
% instruction deliberately does not.
\section{How the Company Holds Its Own Money}

A funder reading a small business proposal reasonably asks how that business
holds the cash it already has. This section answers it in one page. It is
included in a packet sent to program officers because the answer is short, it is
conservative, and it is the kind of thing that is easier to state before it is
asked than after.

\subsection{Why a ladder rather than a sweep}

The company's operating horizon for the coming period is nine months, matching
the duration of a Small Business Innovation Research Phase I award should one be
made \cite{nihsbirseed}. A single sweep balance earns a rate but expresses no
view about when money is needed. A ladder expresses exactly that view: each
maturity is placed near a point at which the company has a known cash need, so no
position has to be sold before maturity to meet one.

The ladder is also a governance instrument. Each maturity is a decision point at
which the horizon is re-read rather than letting an automatic roll make the
decision silently, which is why auto-reinvestment is switched off on every rung.

\begin{appfloat}[!tb]
\begin{appfig}[d2-type / grid beside a measure strip on one scale]
% Left panel: four claim rows on a 0.76cm pitch.  Right panel: a twelve-month
% rule with four rung markers and a shaded nine-month band.  A 10mm corridor
% separates them and no element spans it.
\foreach \x/\w/\lab in {0/42/{Claim on the reserve}, 3.30/18/{Timing}, 4.90/22/{Rung that meets it}}{
  \node[d2cellh,minimum width=\w mm,text width=\w mm] at (\x,0) {\lab};}
\foreach \y/\c/\t/\r in {%
  -0.76/{Operating burn, quarter 1}/{Month 3}/{Rung A},
  -1.52/{Phase I milestones 1 to 3}/{Month 6}/{Rung B},
  -2.28/{Phase I milestones 4 and 5}/{Month 9}/{Rung C},
  -3.04/{Contingency, unclaimed}/{Month 12}/{Rung D}}{
  \node[d2celll,minimum width=42mm,text width=42mm] at (0,\y) {\c};
  \node[d2cell,minimum width=18mm,text width=18mm] at (3.30,\y) {\t};
  \node[d2cellk,minimum width=22mm,text width=22mm] at (4.90,\y) {\r};}
\ptitle{-1.70}{0.62}{Four claims, four maturities}
% Right panel: shifted scope, 0.40cm per month.
\begin{scope}[shift={(7.60,-2.30)}]
  \node[mmband,minimum width=3.6cm,minimum height=0.5cm,anchor=south west] at (0,-0.25) {};
  \axisx{0}{4.9}{0}
  \foreach \m/\lab in {0/0, 1.2/3, 2.4/6, 3.6/9, 4.8/12}{
    \draw[fundgray,line width=0.3pt] (\m,0) -- (\m,-0.12);
    \node[font=\tiny\sffamily,text=fundgray,anchor=north] at (\m,-0.16) {\lab};}
  \node[font=\tiny\sffamily,text=fundgray,anchor=north] at (2.4,-0.52) {months from settlement};
  \foreach \m/\h/\lab in {1.2/0.62/{A, 20 pct}, 2.4/1.30/{B, 20 pct},
                          3.6/0.62/{C, 20 pct}, 4.8/1.30/{D, 20 pct}}{
    \draw[fundmid,line width=0.7pt] (\m,0.04) -- (\m,\h-0.16);
    \node[d2cellk,minimum width=15mm,text width=13mm] at (\m,\h) {\lab};}
  \ptitle{0}{2.10}{The same claims on a twelve-month rule}
  \pnote{0}{-0.90}{Shaded band: the nine-month Phase I horizon.}
  \legkey{0}{-1.25}{fundpale}{Liquid sleeve, 15 pct, no maturity}
  \legkey{2.35}{-1.25}{fundgrayl}{Sweep residual, 5 pct}
\end{scope}
\end{appfig}
\vspace{-0.60cm}
\figcaption{Figure 3. The corporate reserve as four maturities and one liquid sleeve,\\
placed against the dated claims each rung is cut to meet, at one scale.}
\end{appfloat}

\subsection{The six lines}

Sizes are expressed as a share of the reserve rather than in dollars, so the
instruction stays correct whatever the balance is on the day it is approved. No
dollar amount is stated here for the same reason.

\begin{apptable}
\begin{tabularx}{\textwidth}{>{\raggedright\arraybackslash}p{2.2cm} >{\raggedright\arraybackslash}p{3.0cm} >{\raggedright\arraybackslash}p{1.5cm} >{\raggedright\arraybackslash}p{2.6cm} >{\raggedright\arraybackslash}X}
\headrow \thead{Line} & \thead{Instrument} & \thead{Share} & \thead{Order type} & \thead{Maturity or horizon}\\
\midrule
Rung A & U.S. Treasury bill & 20 pct & Non-competitive at auction & About 3 months \\
Rung B & U.S. Treasury bill & 20 pct & Non-competitive at auction & About 6 months \\
Rung C & U.S. Treasury bill & 20 pct & Non-competitive at auction & About 9 months \\
Rung D & U.S. Treasury note or bill & 20 pct & Secondary market & About 12 months \\
Liquid sleeve & Short-duration Treasury ETF & 15 pct & Limit, day, never market & No maturity; sized for an unexpected need \\
Residual & Government money market sweep & 5 pct & Sweep & Overnight \\
\end{tabularx}
\end{apptable}
\vspace{-0.60cm}
\tabcap{Table 4. The six lines of the reserve, each with its instrument, its\\
share, the order type it is entered under, and the horizon it is cut to.}

Rungs A through D are direct obligations with a stated maturity date and a CUSIP
on each rung, not a fund wrapper \cite{treasuryauctions2026}. The distinction is
the whole point: a fund has a price and a direct obligation has a maturity, and
the ladder exists to own maturities.

\subsection{What is excluded, and why}

No margin, no options, no securities lending, no prime money market fund, no
corporate credit, and no equity. A research company's operating reserve should
carry no financing risk, no optionality, and no counterparty risk it was not paid
to take. Award funds, if any arrive, are segregated and never mixed into the
reserve, which is the same separation of federal from non-federal funds the site
package sets out \cite{movein2026}.

Nothing in this section is investment advice or a recommendation to any other
person, and nothing in it has been placed. It is the sole member's direction for
the company's own reserve.
